\begin{abstract}
\textcolor{black}{\textit{Agents}, language model-based systems capable of reasoning, planning, and acting are widely adopted in real-world tasks, yet how their performance changes as these systems scale across key dimensions remains underexplored. We introduce quantitative \textit{scaling principles} for agent systems as a predictive model, capturing how performance varies with coordination, model capability, and measurable system and task factors. Across 260 configurations spanning six agentic benchmarks, five canonical architectures (Single-Agent and four Multi-Agent: Independent, Centralized, Decentralized, Hybrid), and three LLM families, we perform controlled evaluations standardizing tools, prompts, and compute to isolate architectural effects. The resulting model achieves a cross-validated $R^2{=}0.373$ across all six benchmarks ($R^2{=}0.413$ with a task-grounded capability metric). We identify a robust capability-saturation effect and additional patterns: \textbf{(1)} a coordination yields diminishing returns once single-agent baselines exceed certain performance; \textbf{(2)} tool-heavy tasks appear to incur multi-agent overhead; and \textbf{(3)} architectures without centralized verification tend to propagate errors more than those with centralized coordination. Relative performance change compared to single-agent baseline ranges from ${+}80.8\%$ on decomposable financial reasoning to ${-}70.0\%$ on sequential planning, demonstrating that architecture-task alignment determines collaborative success. The framework identifies the best-performing architecture for 87\% of held-out configurations and shows consistent relative architecture preferences on unseen frontier models. Agent effectiveness depends on alignment between coordination and task structure, and that mismatched coordination degrades the performance.}
\end{abstract}